 Este capítulo describe la estrategia técnica y operativa para el despliegue del sistema DOMU en un entorno
      productivo, asegurando su disponibilidad y correcta adopción por parte de la comunidad.
    5
    6 \section{Arquitectura de Despliegue}
    7 Para la puesta en marcha se ha diseñado una arquitectura basada en contenedores, lo que garantiza la portabilidad
      facilita el escalado horizontal.
    8
    9 \begin{itemize}
   10     \item \textbf{Servidor de Aplicaciones:} El backend Java se empaqueta en una imagen Docker basada en
      \texttt{eclipse-temurin:21-alpine}, optimizada para bajo consumo de recursos. Se despliega en un servicio de
      orquestación (como AWS ECS o un VPS con Docker Compose).
   11     \item \textbf{Base de Datos:} Se utiliza una instancia gestionada de MySQL 8.0 en la nube (ej. AWS RDS o Googl
      Cloud SQL) para garantizar copias de seguridad automáticas y alta disponibilidad.
   12     \item \textbf{Frontend y Móvil:} La aplicación web se distribuye a través de una red de entrega de contenido
      (CDN) como Vercel o Netlify. La aplicación móvil se genera mediante los servicios de \textit{build} de Expo (EAS)
      para su distribución en tiendas de aplicaciones o como APK para instalación directa en dispositivos Android.
   13     \item \textbf{Almacenamiento:} Los archivos estáticos (evidencias, boletas, fotos de perfil) se alojan en un
      \textit{bucket} de Google Cloud Storage, configurado con políticas de acceso privado y URLs firmadas temporalmente
   14 \end{itemize}
   15
   16 \section{Plan de Capacitación}
   17 La adopción del sistema requiere nivelar las competencias digitales de los distintos usuarios.
   18
   19 \subsection{Administración y Comité}
   20 Se realiza una sesión de inducción de 2 horas enfocada en la configuración inicial del edificio: carga de la nómin
      de residentes, configuración de alícuotas y definición de áreas comunes. Se entrega un "Manual de Gestión" en
      formato digital.
   21
   22 \subsection{Personal de Conserjería}
   23 Capacitación práctica en terreno (30 minutos) centrada exclusivamente en el uso del módulo de "Control de Accesos"
      la operación de la pistola lectora de QR. Se prioriza la rapidez operativa para no entorpecer el flujo de ingreso.
   24
   25 \subsection{Residentes}
   26 Se distribuye un video tutorial de 3 minutos vía WhatsApp y correo electrónico, explicando cómo descargar la app,
      recuperar la contraseña y registrar una visita o reserva.
   27
   28 \section{Estrategia de Lanzamiento}
   29 Se recomienda una estrategia de \textbf{Marcha Blanca} de dos semanas.
   30 \begin{enumerate}
   31     \item \textbf{Semana 1 (Paralelo):} El sistema DOMU funciona simultáneamente con el libro de visitas manual. N
      se bloquean accesos ni se cobran multas. Objetivo: Poblar la base de datos y habituar al personal.
   32     \item \textbf{Semana 2 (Exclusividad):} Se elimina el registro en papel para visitas y encomiendas. Se habilit
      el módulo de reservas.
   33     \item \textbf{Lanzamiento Oficial:} Se habilita el módulo financiero para el siguiente ciclo de facturación.
   34 \end{enumerate}